
% Kapitola 4: LLM ako podpora správy metadát
% Stav: finálna verzia

\section{LLM ako podpora správy metadát}

Predchádzajúca kapitola ukázala, že pri správe publikačných metadát sa opakovane objavujú situácie, ktoré nemožno uspokojivo vyriešiť len jednoduchou syntaktickou validáciou. Ide najmä o prípady, kde treba interpretovať voľný text, porovnávať variantné zápisy, rozhodovať medzi viacerými kandidátmi alebo prevádzať nejednotné vstupy do kontrolovanej štruktúry. Práve v tomto priestore sa otvára miesto pre využitie veľkých jazykových modelov ako asistenčnej vrstvy nad tradičnými heuristikami a validačnými pravidlami.

Táto kapitola preto nepojednáva o LLM ako o všeobecnom fenoméne, ale výlučne v rozsahu potrebnom pre túto prácu. Dôraz sa kladie na tie vlastnosti modelov, ktoré sú priamo relevantné pre extrakciu, normalizáciu a kontrolu bibliografických údajov, na návrh ich vstupov a výstupov a na podmienky, za ktorých zostáva ich nasadenie bezpečné, auditovateľné a metodicky obhájiteľné.

\subsection{Stručný princíp LLM}

Veľké jazykové modely sú generatívne modely trénované na rozsiahlych textových korpusoch, ktoré dokážu na základe vstupného kontextu predikovať a vytvárať ďalší text. Pre potreby tejto práce je dôležité najmä to, že model dokáže pracovať s neštruktúrovaným alebo len čiastočne štruktúrovaným vstupom a prevádzať ho do podoby, ktorá lepšie zodpovedá požiadavkám následného strojového spracovania \cite{AlammarGrootendorst2024,Foster2019}. V praxi to znamená, že model môže pomôcť tam, kde sa v jednom poli mieša faktická informácia, prirodzený jazyk a nejednotný zápis.

Zároveň však ide o model pravdepodobnostný, nie deterministický. To znamená, že ani pri rovnakom zadaní nemožno bez ďalších opatrení automaticky predpokladať úplne totožný alebo bezchybný výstup. Yanev zdôrazňuje, že pri budovaní aplikácií nad LLM je potrebné navrhovať ich použitie tak, aby sa ich výstupy dali validovať, obmedziť a spoľahlivo začleniť do širšej aplikačnej logiky \cite{Yanev2024}. V kontexte tejto práce je preto podstatné chápať LLM nie ako náhradu evidenčného systému, ale ako komponent určený na podporu konkrétnych kurátorských úloh.

\subsection{Roly LLM v metadátovom workflow}

V metadátovom workflow možno jazykový model chápať minimálne v troch odlišných rolách. Prvou je rola normalizátora, teda nástroja, ktorý z nejednotného textového vstupu vytvára kandidátnu normalizovanú hodnotu alebo štruktúrovaný návrh zodpovedajúci požiadavkám systému. Druhou je rola asistenta rozhodovania, v ktorej model pomáha používateľovi pri výbere medzi viacerými kandidátmi alebo pri interpretácii nejednoznačného poľa. Tretia rola je rola sekundárneho kontrolóra konzistencie, keď model upozorňuje na možné rozpory medzi poliami alebo na neúplnosť záznamu \cite{Mikulecky2025}.

Tieto roly sa navzájom odlišujú mierou rizika aj typom očakávaného výstupu. Pri normalizácii je cieľom najmä transformácia voľného textu do kontrolovanej podoby. Pri asistencii rozhodovania je dôležitejšia argumentovaná voľba medzi alternatívami. Pri kontrole konzistencie ide zasa o schopnosť modelu zachytiť, že niektoré polia si navzájom odporujú alebo že záznam nezodpovedá očakávanej logike. Mikulecký ukazuje, že práve úlohy viazané na extrakciu a transformáciu polí patria medzi najsľubnejšie oblasti nasadenia AI v evidencii vedeckej publikačnej činnosti \cite{Mikulecky2025}. Pre túto prácu je však rozhodujúce, že vo všetkých troch rolách musí model fungovať ako asistent v rámci kontrolovaného workflow, nie ako autonómny vykonávateľ zápisov do databázy.

\subsection{Štruktúrované výstupy a prompt engineering}

Pri spracovaní bibliografických metadát nestačí, aby model vytvoril jazykovo presvedčivú odpoveď. Výstup musí byť navrhnutý tak, aby sa dal spoľahlivo parsovať, validovať a uložiť do aplikácie. Preto je v tejto doméne dôležitejšia schopnosť generovať štruktúrované odpovede než schopnosť formulovať voľný text. OpenAI aj Ollama dnes explicitne podporujú structured outputs, teda režim, v ktorom sa odpoveď modelu podriaďuje zadanej JSON schéme \cite{OpenAIStructuredOutputs,OllamaStructuredOutputs}. Takýto prístup významne znižuje rozdiel medzi jazykovým modelom a bežnou aplikačnou logikou, pretože výstup modelu možno bezprostredne podrobiť formálnej validácii.

S tým úzko súvisí aj prompt engineering, teda návrh vstupov tak, aby bol model vedený k čo najstabilnejšiemu a najrelevantnejšiemu správaniu. V tejto práci však nejde o prompt engineering v širokom zmysle, ale o jeho úzko praktickú podobu zameranú na extrakciu polí, dopĺňanie kandidátnych hodnôt a kontrolu konzistencie. Yanev aj Alammar s Grootendorstom sa zhodujú, že pri aplikáciách orientovaných na dáta je potrebné zadanie formulovať presne, minimalizovať nejednoznačnosť a explicitne špecifikovať požadovaný výstupný formát \cite{Yanev2024,AlammarGrootendorst2024}.

\subsubsection{JSON schéma a function calling}

Štruktúrovaný výstup založený na JSON schéme má v tejto práci zásadný význam. Umožňuje totiž presne určiť, aké polia sa od modelu očakávajú, aký majú mať typ, ktoré z nich sú povinné a aké obmedzenia sa na ne vzťahujú. OpenAI rozlišuje structured outputs pre odpovede priamo v JSON schéme a function calling, pri ktorom model vytvára argumenty pre vopred definovaný nástroj alebo funkciu \cite{OpenAIStructuredOutputs,OpenAIFunctionCalling}. Ollama podporuje obdobný princíp tým, že modelu možno pri chat rozhraní odovzdať JSON schému cez parameter \texttt{format} \cite{OllamaStructuredOutputs}.

Pre navrhované riešenie je tento prístup dôležitý preto, že oddeľuje jazykovú interpretáciu vstupu od následnej systémovej akcie. Model môže napríklad navrhnúť štruktúrovaný objekt s menom autora, kandidátnou afiliáciou alebo normalizovaným dátumom, no až ďalšia vrstva aplikácie rozhodne, či je tento návrh formálne validný, či je v zhode s autoritami a či bude predložený používateľovi na schválenie. Function calling je v tomto kontexte užitočný najmä preto, že model nevracia iba text, ale argumenty pre presne vymedzenú operáciu, čo zvyšuje kontrolovateľnosť celého procesu.

\subsubsection{Prompt engineering}

Prompt engineering má v tejto doméne zmysel predovšetkým ako nástroj obmedzenia variability modelových výstupov. Oficiálna dokumentácia OpenAI zdôrazňuje potrebu klásť inštrukcie jednoznačne, jasne oddeľovať inštrukcie od vstupných dát, byť konkrétny v požadovanom výstupe a pri zložitejších úlohách používať aj few-shot príklady \cite{OpenAIPromptGuide}. Pri metadátových úlohách je dôležité, aby model presne vedel, ktoré časti vstupu má chápať ako dáta na analýzu a ktoré ako pravidlá spracovania.

Z praktického hľadiska to znamená, že systémový prompt má vymedziť úlohu modelu a hranice jeho správania, zatiaľ čo používateľský vstup alebo importovaný text predstavuje materiál na spracovanie. Few-shot príklady sú užitočné najmä pri poliach s opakujúcou sa štruktúrou, napríklad pri extrakcii dátumov z redakčných poznámok alebo pri mapovaní autorov na očakávanú výstupnú štruktúru. Cieľom nie je vytvoriť univerzálny prompt pre všetky situácie, ale navrhnúť menší počet stabilných promptových šablón pre konkrétne typy úloh.

\subsection{Human-in-the-loop a riziká}

Použitie LLM pri správe metadát má zmysel iba vtedy, ak zostáva človek súčasťou rozhodovacieho procesu. Human-in-the-loop v tomto kontexte neznamená iba formálnu možnosť zásahu, ale vedomý návrhový princíp, podľa ktorého model pripravuje návrh a používateľ nesie zodpovednosť za jeho potvrdenie alebo odmietnutie. Tento princíp je dôležitý najmä tam, kde model pracuje s nejednoznačnými údajmi, s variantnými zápismi autorov alebo s dopĺňaním hodnôt, ktoré nemožno deterministicky overiť \cite{Mikulecky2025}.

Dôvodom nie je len metodická opatrnosť, ale aj bezpečnostné riziká. OWASP upozorňuje, že LLM aplikácie sú zraniteľné voči prompt injection, pri ktorej útočník alebo škodlivý vstup manipuluje správanie modelu prostredníctvom prirodzeného jazyka \cite{OWASPPromptInjection}. Zároveň ide o systémy náchylné na halucinácie, teda na produkciu presvedčivo znejúcich, no vecne nesprávnych hodnôt. V širšom bezpečnostnom rámci OWASP patrí medzi relevantné riziká aj únik citlivých informácií, neprimeraná dôvera v modelové výstupy, nedostatočná kontrola nástrojov a závislosť od dodávateľského reťazca modelov a knižníc \cite{OWASPTop10LLM2025}. V doméne bibliografických metadát síce zvyčajne nejde o vysoko citlivé osobné údaje, no aj nesprávne alebo nekontrolovane prijaté výstupy môžu poškodiť kvalitu evidencie a dôveryhodnosť systému.

Preto sa ako vhodný model javí taký návrh, v ktorom LLM navrhuje, no nekoná sám. Výstup modelu má byť validovaný, auditovaný a v rizikových prípadoch predkladaný používateľovi ako kandidátna hodnota. Takýto prístup znižuje nielen pravdepodobnosť chybného zápisu, ale aj prevádzkové riziká spojené s tým, že model by mal priamy a nekontrolovaný dosah na interné údaje.

\subsection{Lokálne vs cloud nasadenie}

Pri integrácii LLM do informačného systému treba rozhodnúť aj o spôsobe nasadenia. V zásade ide o voľbu medzi lokálnym režimom, v ktorom model beží v prostredí organizácie alebo vývojára, a cloudovým režimom, v ktorom sa model používa prostredníctvom vzdialeného API. Ollama reprezentuje typický príklad lokálneho prístupu, pri ktorom je model obsluhovaný cez lokálne rozhranie a možno ho začleniť do aplikácie bez závislosti od vzdialeného poskytovateľa \cite{OllamaDocs}. Na opačnej strane stoja cloudové služby s OpenAI-kompatibilným rozhraním, pri ktorých je výhodou jednoduchšia integrácia a často aj vyšší výkon alebo širší výber modelov; príkladom je dokumentovaná kompatibilita Groq API s klientskymi knižnicami OpenAI \cite{GroqOpenAICompatibility}.

Rozhodovanie medzi týmito režimami nemožno zužovať na jediný parameter. Lokálne nasadenie znižuje závislosť od externého poskytovateľa a zjednodušuje kontrolu nad tokom dát. Cloudové nasadenie zasa často ponúka nižšie bariéry pri škálovaní, menšiu prevádzkovú záťaž a prístup k modelom, ktoré by lokálne nebolo praktické prevádzkovať. Yanev pri budovaní aplikácií nad LLM upozorňuje, že architektúra systému musí zohľadniť nielen kvalitu modelu, ale aj latenciu, spoľahlivosť rozhrania, cenu používania a spôsob validácie odpovedí \cite{Yanev2024}. V tejto práci preto nejde o normatívny výber jednej možnosti, ale o návrhový rámec, v ktorom môže byť LLM vrstva vymeniteľná podľa prevádzkových potrieb.

\subsection{Vyhodnocovanie úspešnosti}

Ak má byť použitie LLM v práci metodicky obhájiteľné, nestačí konštatovať, že model „pomáha“. Potrebné je určiť, v ktorých úlohách sa jeho prínos meria, aký je referenčný stav a akým spôsobom sa bude vyhodnocovať správnosť alebo užitočnosť výstupu. V kontexte tejto práce je prirodzené hodnotiť úspešnosť na úrovni jednotlivých polí, nie iba na úrovni celého záznamu. Osobitne možno posudzovať napríklad správnosť identifikácie interného autora, normalizácie názvu časopisu, extrakcie dátumu alebo návrhu deduplikačného rozhodnutia.

Mikulecký pri hodnotení AI podpory evidencie publikačných výstupov pracuje s porovnaním voči očakávanému výsledku a ukazuje, že relevantné je sledovať nielen všeobecnú úspešnosť, ale aj typy chýb viazané na konkrétne polia \cite{Mikulecky2025}. V oblasti párovania a entity resolution sa tradične používajú metriky typu precision a recall, prípadne ich kombinácie, pričom Christen zdôrazňuje význam kvalitne pripraveného referenčného súboru, voči ktorému sa výsledok porovnáva \cite{Christen2012}. Pre túto prácu z toho vyplýva potreba gold standardu, teda manuálne overenej vzorky záznamov, na ktorej bude možné overiť prínos navrhovaného riešenia.

Pri úlohách, ktoré obsahujú určitú mieru interpretačnej nejednoznačnosti, je vhodné sledovať aj zhodu medzi ľudskými anotátormi. Cohenova kappa je klasická miera zhody pre nominálne kategórie a umožňuje odlíšiť skutočnú zhodu od zhody, ktorá by mohla vzniknúť náhodne \cite{Cohen1960}. V kontexte tejto práce má takáto miera význam najmä pri posudzovaní toho, či je určitý záznam duplicitný, či daný autor patrí medzi interných autorov alebo či je navrhnutá normalizácia ešte prijateľná z pohľadu knihovníckej praxe. Vyhodnocovanie úspešnosti preto nebude iba technickým doplnkom implementácie, ale nevyhnutnou súčasťou argumentácie, že integrácia LLM do workflow prináša merateľný prínos.
